\section{Introduction}

\subsection{The problem}

Every crewed lunar architecture that has flown, and every one currently
proposed, splits the vehicle in two. A descent stage carries the landing gear,
the throttleable engine and the surface consumables, and it stays where it
lands. Only an ascent stage returns to orbit. The reason is arithmetic: mass
that does not have to be re-accelerated to lunar orbital speed does not have to
be carried from Earth in the first place. Apollo used this split six times, and
the Artemis Human Landing System requirements inherit it \cite{nasa_artemis}.

The consequence is that the return leg is a rendezvous problem, and a harder one
than the launch. The ascent stage flies on a single engine with no aerodynamic
stabilisation, no throttle authority, a centre of mass that migrates as the
tanks drain, and an inertia tensor with significant off-diagonal terms. Its
guidance can steer the thrust vector and choose when to stop, but it cannot trim
the acceleration profile. Whatever mismatch remains between predicted and actual
performance appears directly as a velocity error at cut-off. From orbit, an
ascent stage is completely described by the dispersed state it delivers, and
everything after that point is somebody else's propellant budget.

This report follows that budget from the parking orbit to the moment a robotic
arm closes on a grapple fixture.

\subsection{Scope and mission definition}

Table~\ref{tab:mission} defines the case study. Both vehicles are on circular,
equatorial, coplanar orbits about the Moon; the chaser is a lunar module (LM) at
$100\unit{km}$ altitude, the target a mothership (MS) at $400\unit{km}$, leading
by $\phi_0 = 10^\circ$ at the epoch.

\begin{table}[htbp]
\centering
\caption{Mission definition and physical constants. Values are held in a single
struct, \texttt{config/mission\_constants.m}; no physical number appears as a
literal anywhere else in the code base.}
\label{tab:mission}
\small
\begin{tabular}{llll}
\toprule
Quantity & Symbol & Value & Unit \\
\midrule
Lunar gravitational parameter & $\mu_M$    & $4902.8$   & km$^3$/s$^2$ \\
Lunar mean radius             & $R_M$      & $1737.4$   & km \\
Lunar oblateness              & $J_2$      & $2.033\times10^{-4}$ & -- \\
Earth gravitational parameter & $\mu_E$    & $3.986\times10^{5}$  & km$^3$/s$^2$ \\
Mean Earth--Moon distance     & $d_{EM}$   & $384\,400$ & km \\
Chaser orbit radius           & $R_1$      & $1837.4$   & km \\
Target orbit radius           & $R_2$      & $2137.4$   & km \\
Initial target lead           & $\phi_0$   & $10$       & deg \\
CR3BP mass parameter          & $\mu$      & $0.01215$  & -- \\
Random seed                   & --         & $42$       & -- \\
\bottomrule
\end{tabular}
\end{table}

The reference frame for the orbital work is Moon-centred inertial (MCI),
equatorial, with $+z$ along the lunar north pole and motion counter-clockwise in
the $xy$-plane. At $t = 0$ the chaser is on the $+x$ axis. Proximity operations
use a local-vertical/local-horizontal frame attached to the target, defined
carefully in Section~\ref{sec:proximity} because its handedness is the single
most error-prone convention in the whole study.

\subsection{What is modelled, and what is not}

This is an engineering-faithful simulation of a restricted model. It is not a
digital twin of any flown or planned mission, and it is worth being explicit
about the boundary before presenting results.

\emph{Modelled:} two-body dynamics with tangential impulsive manoeuvres; orbital
phasing; linear and exact nonlinear relative motion; Cowell propagation with
lunar $J_2$ and Earth third-body perturbation; the Earth--Moon circular
restricted three-body problem; the kinematics and kineto-statics of a
five-revolute manipulator on a controlled base.

\emph{Not modelled:} the lunar gravity field beyond $J_2$, in particular the
mass concentrations mapped by GRAIL \cite{zuber_grail} and the $C_{22}$
ellipticity term; solar radiation pressure and eclipse cycling; finite-burn
losses, propellant slosh and structural flexibility; orbital inclination and
nodal regression; a near-rectilinear halo orbit in place of the circular target
orbit; and the free-floating multi-body dynamics of the arm, whose inertia
matrix and Coriolis terms are computed but never integrated. Section~\ref{sec:limits}
returns to each of these and states what it would change.

\subsection{Numerical approach and verification philosophy}

Everything is implemented from first principles in MATLAB. The only solvers used
are \texttt{ode45} for propagation and \texttt{fminsearch} for the three-body
shooting problem; no MathWorks toolbox beyond base MATLAB is required, and no
external ephemeris is read.

The organising principle of the work is that a simulation is worth exactly as
much as its verification. Each part therefore closes with a verification
subsection stating what was compared against what. Three kinds of check are used
throughout:

\begin{enumerate}[label=(\roman*)]
\item \textbf{Closed form.} Where an analytical solution exists it is
      implemented independently and differenced against the numerical one. The
      Kepler propagation check of Section~\ref{sec:transfer:verif} is the
      strongest example.
\item \textbf{Invariants.} Quantities the dynamics must conserve are monitored:
      specific energy and angular momentum on each Keplerian leg, and the Jacobi
      constant along the three-body arc.
\item \textbf{Independent implementation.} Selected results are compared against
      a separate implementation of the same mission geometry. Where the two
      disagree, Section~\ref{sec:discussion:verif} says by how much and why,
      rather than tuning one towards the other.
\end{enumerate}

A dedicated script, \texttt{tests/audit\_reference.m}, performs all of these
automatically at the end of every build and writes a pass/warn/fail table. The
results quoted throughout this report are the ones that script validated.
